\documentclass[12pt, english]{article}
%%%%%%%%%%%%%%%%%%%%%%%%%%%%%%%%%%%%%%%%%%%%%%%%%%%%%%%%

% margin %%%%%%%%%%%%%%%%%%%%%%
\usepackage{geometry}
 \geometry{
 a4paper,
 left=1in,
 top=0.9in,
 right=1in,
 bottom=2in,
 }
%%%%%%%%%%%%%%%%%%%%%%%%%%%%%%%%%%%%%%%%%%%%%%%%%%%%%%%%

% page no. %%%%%%%%%%%%%%%%%%%%%%
\usepackage{fancyhdr}
\pagestyle{fancy}
\fancyhf{}
\fancyfoot[C]{\textit{\thepage}}
\setlength{\headheight}{65pt}
%%%%%%%%%%%%%%%%%%%%%%%%%%%%%%%%%%%%%%%%%%%%%%%%%%%%%%%%

% idk-needed %%%%%%%%%%%%%%%%%%%%%%
\usepackage{amsmath,amsthm,amssymb,amsfonts, enumitem, color, comment, graphicx, environ, lipsum, accsupp}
%%%%%%%%%%%%%%%%%%%%%%%%%%%%%%%%%%%%%%%%%%%%%%%%%%%%%%%%

% code and color %%%%%%%%%%%%%%%%%%%%%%
\usepackage{spverbatim}
\usepackage{xcolor}
\usepackage{listings}
\definecolor{bluegrey}{RGB}{0, 0, 200}
\lstset{
        framexleftmargin=0pt,
        frame=shadowbox,
        rulesepcolor=\color{gray},
        xleftmargin=30pt,
        xrightmargin=30pt,
        breaklines=true,
        basicstyle=\ttfamily,
        showstringspaces=false,
        }
%%%%%%%%%%%%%%%%%%%%%%%%%%%%%%%%%%%%%%%%%%%%%%%%%%%%%%%%

% heading %%%%%%%%%%%%%%%%%%%%%%
\renewcommand\thesection{\color{blue}\arabic{section}}
\newcommand{\mysection}[2]{\setcounter{section}{#1}\addtocounter{section}{-1}\section{#2}}
%%%%%%%%%%%%%%%%%%%%%%%%%%%%%%%%%%%%%%%%%%%%%%%%%%%%%%%%

% both side indent %%%%%%%%%%%%%%%%%%%%%%
\newenvironment{myindent}
{\par\leftskip25pt\relax\rightskip25pt\relax}
{\par\leftskip0cm\relax\rightskip0cm\relax}
%%%%%%%%%%%%%%%%%%%%%%%%%%%%%%%%%%%%%%%%%%%%%%%%%%%%%%%%

% type ^ and \ %%%%%%%%%%%%%%%%%%%%%%
\newcommand{\power}{$^\wedge$}
\newcommand{\backs}{\textbackslash} 
%%%%%%%%%%%%%%%%%%%%%%%%%%%%%%%%%%%%%%%%%%%%%%%%%%%%%%%%

% document info %%%%%%%%%%%%%%%%%%%%%%
\lhead{Arav Sarma \\ 15193845}  %replace with your name
\rhead{Homework \textcolor{blue}{4} \\ 2025 Speech Processing}
%%%%%%%%%%%%%%%%%%%%%%%%%%%%%%%%%%%%%%%%%%%%%%%%%%%%%%%%

% begin %%%%%%%%%%%%%%%%%%%%%%
\begin{document}
%wwwwwwwwwwwwwwwwwwwwwwwwwwwwwwwwwwwwwwwwwwwwwwwwwwwwwwwww

% section _______________________________________________
\mysection{4}{Files}
\vspace{10pt}
% -__-__-__-__-__-__-__-__-__-__-__-__-__-__-__-__-__-__-

% //////////////////////////////////////////////////////
\subsection*{\emph{1.} Path to sound files} 
\vspace{10pt}

\begin{itemize}
    \item Full path name: \textbf{D:/Desktop/spTest/sounds/tone1.wav}
    \item Relative path from spTest: \textbf{spTest/sounds/tone1.wav}
    \item \textcolor{blue}{Relative path from spTest: \textbf{sounds/tone1.wav}}
    \item Relative path from scripts: \textbf{../sounds/tone1.wav}
\end{itemize} 

% //////////////////////////////////////////////////////

% \\\\\\\\\\\\\\\\\\\\\\\\\\\\\\\\\\\\\\\\\\\\\\\\\\\\\\\\\\
\subsection*{\emph{2.} Vector using fileNames\$\#}

\begin{lstlisting}
files$# = fileNames$# ("../sounds/*.wav")
writeInfo: files$#
\end{lstlisting}

% \\\\\\\\\\\\\\\\\\\\\\\\\\\\\\\\\\\\\\\\\\\\\\\\\\\\\\\\\\

% \\\\\\\\\\\\\\\\\\\\\\\\\\\\\\\\\\\\\\\\\\\\\\\\\\\\\\\\\\
\subsection*{\emph{3.} Retrieve vector items using size()}

\begin{lstlisting}
writeInfo: ""
files$# = fileNames$# ("../sounds/*.wav")
num = size(files$#)
for i to num
	appendInfoLine: files$#[i]
endfor
\end{lstlisting}

% \\\\\\\\\\\\\\\\\\\\\\\\\\\\\\\\\\\\\\\\\\\\\\\\\\\\\\\\\\

% \\\\\\\\\\\\\\\\\\\\\\\\\\\\\\\\\\\\\\\\\\\\\\\\\\\\\\\\\\
\subsection*{\emph{4.} Get duration of sound files}

\begin{myindent}
    \texttt{sound = Read from file: "tone1.wav"} does not work as the file is not in the same directory as the script. Without specifying the path, the code looks for the file \texttt{tone1.wav} in the same directory the script is in.
\end{myindent}
\vspace{10pt}

\begin{lstlisting}
writeInfo: ""
files$# = fileNames$# ("../sounds/*.wav")
for i to size(files$#)
	Read from file: "../sounds/" + files$#[i]
	duration = Get total duration
	appendInfoLine: duration
endfor
\end{lstlisting}

% \\\\\\\\\\\\\\\\\\\\\\\\\\\\\\\\\\\\\\\\\\\\\\\\\\\\\\\\\\

% \\\\\\\\\\\\\\\\\\\\\\\\\\\\\\\\\\\\\\\\\\\\\\\\\\\\\\\\\\
\subsection*{\emph{5.} Save fragments of sound files}

\begin{lstlisting}
writeInfo: ""
files$# = fileNames$# ("../sounds/*.wav")
newPathName$ = "../fragments"
for i to size(files$#)
	fragments$ = files$#[i] - ".wav" + "_short.wav"
	Read from file: "../sounds/" + files$#[i]
	Extract part: 0, 0.1, "rectangular", 1, "no"
	Save as WAV file: newPathName$ + "/" + fragments$
endfor
\end{lstlisting}

\vspace{10pt}
% \\\\\\\\\\\\\\\\\\\\\\\\\\\\\\\\\\\\\\\\\\\\\\\\\\\\\\\\\\

% ><><><><><><><><><><><><><><><><><><><><><><><><><><><><><
\subsection*{\emph{6.} The continua of your experiment}
\vspace{10pt}

\begin{itemize}
    \item \textbf{General research question:} \par
    \textcolor{gray}{Location of phoneme boundary as a function of orthographic context (and exopure/familiarity with orthography of certain language(s))} \par
    \textbf{Location of phoneme boundary as a function of lexical context (word-nonword)}
    \item \textbf{\textcolor{gray}{Addition:}} \par
    \textcolor{gray}{L2 lexical knowledge context}
    \item \textbf{Specific research question:} \par
    \textcolor{gray}{My idea is to take both orthography and non-native lexical knowledge as a coupled effect on phonological categorisation. Preliminary thoughts on operationalising this question: looking at phonemic boundary between voiced-plain-aspirated, aspirated-fricative continua for L2 English speakers from Assam (or L1 Assamese speaking). Onset stop contrast on L2 variety is through voicing unlike aspiration in L1 English for Assamese (voicing+aspiration for Boro), Orthograhpic <b>, <d>, <g> are percieved as voiced like other Indian Englishes. This perception gets carried over even when exposed to L1 English. Questions: How much closure devoicing / afterburst duration can be blocked off during phoneme categorisation for a voicing / plain percept (how strong would lexical effect be on a word like `ban' / `pan', and also orthographic; eg: when shown <pan> / <phan or fan>)?} \par
    How does presence/absence of a lexical item effect categorisation of aspirated-breathy laryngeal contrast? Unlike the voicing or aspirated contrasts where boundary can be tested with VOT continua, the aspirated-breathy contrast would require manipulation of cues like turbulence quality and also pitch on following vowel (as the breathy and aspirated categories has been observed to effect pitch across languages; tonogenesis in Rohingya, Punjabi, tone-depressing effect in Southeastern Bantu, adaptation of indic loans with breathy in Thai). The study will look at lexical effect across this contrast for L1 Assamese speakers from Guwahati.       
    \item \textbf{Participants} \par
    Recruited: \char`\~ 5 \\
    Completed: 0 \par
    Participants are going to be tested online using Zoom online meetings, and recruited through author's contacts, and likely by reaching out to Assam/Northeast related online forums like on reddit. The plan is to have non-native English speakers ideally all from Guwahati with same L1, this is to neutralise individual factors like regional variation, exposure to English. A different experiment setup can be by using a qualtrics form along with online meet, so the participants have more autonomy with the test and excess to stimuli on their device while also making it easier to see if they are in a quiet environment. \par
    \item \textbf{Continua}
    Acoustic continua from [bha] to [pha]: 1st context is with following syllable [ŋa] (word to nonword). 2nd context is with following syllable [ki] (nonword to word). \\
    Continua from [dh] to [th]ː 1st context is with nucleus [ue] (word to nonword). 2nd context is with nucleus [io] (nonword to word). \par
\end{itemize}
% ><><><><><><><><><><><><><><><><><><><><><><><><><><><><><

%wwwwwwwwwwwwwwwwwwwwwwwwwwwwwwwwwwwwwwwwwwwwwwwwwwwwwwwww
\end{document}